\section{Considerações finais do exercício}

O esqueleto organiza a pergunta, as referências de apoio e o percurso documental proposto para investigar indicadores de violência letal no Brasil. A contribuição esperada da pesquisa é tornar rastreáveis as escolhas de construção e apresentação dos indicadores e sua mobilização em arenas institucionais, jornalísticas e estatais. Essa expectativa permanece distinta de uma conclusão empírica. A execução do exercício permite verificar a composição do documento, suas referências cruzadas e a correspondência entre fontes bibliográficas e citações, sem substituir coleta, codificação ou interpretação de evidências.

O uso de inteligência artificial neste exercício auxilia a organização do texto e a preparação do projeto LaTeX, mantendo a responsabilidade intelectual do autor. As referências acadêmicas possuem identificadores verificáveis, e a citação direta foi conferida no documento autoral fornecido. A continuidade da pesquisa dependerá da constituição do corpus e da aplicação do protocolo metodológico, sem modificar o recorte consolidado. O documento não registra defesa, banca aprovada, catalogação oficial ou resultados concluídos; sua finalidade é didática e deverá permanecer explícita na revisão e em eventual publicação do fork.
